% Part: incompleteness
% Chapter: theories-computability
% Section: oconsis-ext-of-q-undec

\documentclass[../../../include/open-logic-section]{subfiles}

\begin{document}

\olfileid{inc}{tcp}{oqn}

\olsection{امتدادات $\Th{Q}$ المتسقة-$\omega$ غير قابلة للقرار}

\begin{explain}
كان مدار برهان تمام $\Th{Q}$ بين المجموعات القابلة للتعداد حسابيًا
على أن كل ما يثبت في $\Th{Q}$ من الجمل «صادق» في الأعداد الطبيعية.
ونعيّن الآن، بتعريف ومبرهنة، القدر الذي احتجنا إليه من هذا الصدق،
فنحصل على نتيجة أقوى. ولا حاجة إلى الإطالة في الوقوف عندها، فإننا
سنقويها مرة أخرى بعد قليل. وإنما أوردناها بالدرجة الأولى لبيان
التطور التاريخي: فقد بنى غودل بحثه الأول على الاتساق-$\omega$،
ثم أُثبت بعده بقليل أن الاتساق-$\omega$ يمكن الاستعاضة عنه بالاتساق
المعتاد، فكانت النتيجة أقوى.
\end{explain}

\begin{defn}
نقول إن النظرية $\Th{T}$ متسقة-$\omega$ متى استوفت الشرط الآتي:
لكل جملة $\lexists[x][!A(x)]$، إذا أثبتت $\Th{T}$ جميع الجمل
$\lnot !A(\num 0)$ و$\lnot !A(\num 1)$ و$\lnot !A(\num 2)$،
\dots، امتنع أن تثبت $\Th{T}$ الجملة $\lexists[x][!A(x)]$.
\end{defn}

\begin{thm}
  \ollabel{thm:oconsis-q}
  إذا كانت $\Th{T}$ نظرية متسقة-$\omega$ تشمل $\Th{Q}$،
  فإن $\Th{T}$ غير قابلة للقرار.
\end{thm}

\begin{proof}
احتواء $\Th{T}$ على $\Th{Q}$ يجعل $\Th{T}$ قادرة على تمثيل الدوال
والعلاقات القابلة للحساب. فنكتفي بتعديل البرهان المتقدم:
إذا كان $x\in K$، أثبتت $\Th{T}$، كما سبق،
$\lexists[s][!A_T(\num x,\num x,s)]$. ولإثبات العكس نفرض أن
$\Th{T}$ تثبت $\lexists[s][!A_T(\num x,\num x,s)]$.
يلزم عندئذ انتماء $x$ إلى $K$؛ إذ لو لم ينتم إليها لما كان للآلة
$x$ حساب يتوقف على الدخل $x$. وبما أن $!A_T$ تمثل علاقة كليني $T$،
لأثبتت $\Th{T}$ كلًا من
$\lnot !A_T(\num x,\num x,\num 0)$ و
$\lnot !A_T(\num x,\num x,\num 1)$، \dots. وهذا ينافي كون
$\Th{T}$ متسقة-$\omega$.
\end{proof}

\end{document}
